\section{Introduction}
\label{sec:intro}

Set the temperature to zero and an LLM should give the same answer to the same prompt every time. Evaluation harnesses, regression tests, answer caches and LLM-as-judge pipelines all rely on this assumption. In practice, however, prompts are almost never decoded alone. Serving engines such as \vllm~\citep{kwon2023vllm} batch many requests together to use the GPU efficiently, and which requests share a batch depends on traffic, on the evaluation harness, and on configuration choices nobody records.

{\bf Does the batch a prompt lands in change its greedy output?} Floating-point addition is not associative~\citep{goldberg1991floating}. If a kernel picks a different reduction order or tiling for a different batch shape, the logits of one request can depend on the requests it is co-scheduled with. \citet{he2025nondeterminism} argue that this lack of \emph{batch invariance}, rather than run-to-run randomness from atomics, is the main source of LLM inference nondeterminism. \citet{yuan2025fp32} show that batch size, GPU count and dtype change the accuracy of reasoning models. These works establish that the problem exists. Practitioners still lack a clear picture of its size and shape: how many prompts change, where in the sequence the change begins, which tasks and models are affected, and whether accuracy moves.

{\bf What is missing?} Most reports vary several factors at once. Batch size changes together with batch composition, task type changes together with output length, and model size changes together with numeric format. Run-to-run noise is seldom measured as a separate floor. As a result, common intuitions remain untested, such as ``math diverges more than classification'' or ``larger models are more stable''. The token-level cause of each divergence is also rarely examined.

{\bf Our approach.} We run a controlled, per-prompt study on open models in a single \vllm engine per model. We submit 1,700 byte-identical prompts in static chunks of $B\in\{1,4,16,64\}$, and add three controls: an exact repeat of the $B=64$ schedule, a repeat of the $B=1$ schedule, and a reshuffle that keeps $B=64$ but changes batch-mates (\tabref{tab:conditions}). The prompts come from five tasks that cross math, knowledge and classification with long and short outputs. In particular, MMLU appears both with chain-of-thought (CoT) and as a single-letter answer, which lets us separate task from length. We record the top-2 log-probabilities of every generated token, so that we can inspect the reference model's margin at the exact token where two runs first differ. We evaluate five bf16 models from 0.5B to 3B parameters in two families, plus fp16, fp32 and AWQ variants of \qwen.

{\bf What do we find?} Batching changes far more outputs than expected. For bf16 models, 20--34\% of all prompts and 47--76\% of CoT generations differ between $B=1$ and $B=64$ (\figref{fig:batch}). Already at $B=4$ the effect is at full size, and reshuffling batch-mates at fixed $B=64$ produces 83--100\% of it, while identical schedules agree on all 19,720 repeated generations. The mechanism is simple: about 45\% of first-divergence tokens are exact bf16 ties and at least 91\% are within one rounding step (0.125 nats), a 40$\times$ enrichment over the base rate of ties. Raising precision on fixed weights cuts divergence from 33.5\% (bf16) to 7.3\% (fp16) and 2.8\% (fp32). Mean accuracy stays within $\pm$1.8\pp, with no significant change after Holm correction, but 3--5\% of items flip correctness.

Our contributions are:
\begin{itemize}[leftmargin=*,itemsep=0pt,topsep=0pt]
    \item We measure how often greedy outputs change with batching across ten model configurations and 1,700 prompts, and separate the effect of batch size, batch composition and run-to-run noise. Composition, not size, is the variable that matters, and the engine is fully deterministic for a fixed schedule.
    \item We trace every bf16 divergence to a near-tied top-2 logit and show a monotone precision ladder (bf16 $>$ fp16 $>$ fp32) on identical weights, which identifies rounding of near-ties as the cause.
    \item We show that the apparent stability of classification is entirely an output-length effect (task type n.s., $p=0.58$, once log length is controlled), and that model size between 0.5B and 3B has no consistent effect ($\rho=-0.05$).
    \item We quantify the accuracy impact: mean accuracy is robust, but item-level correctness is not, which matters for regression tests, caching and paired model comparisons.
\end{itemize}

The rest of the paper is organized as follows. \Secref{sec:related} reviews related work, \secref{sec:method} describes our setup, \secref{sec:results} presents results for each hypothesis, \secref{sec:discussion} discusses implications and limitations, and \secref{sec:conclusion} concludes.
